% -----------------------------------------------------------
\section{Transgression, Transgress}
% -----------------------------------------------------------
	\label{TRANSGRESSION}
	
% % ----------------------
% \subsection{See also}
% % ----------------------

% ----------------------
\subsection{1828 American Dictionary of the English Language} % *1828
% ----------------------
	\label{TRANSGRESSION-1828}

	\begin{compactitem}
		\item The act of passing over or beyond any law or rule of moral duty; the violation of a law or known principle of rectitude; breach of command.
		\item[1] Fault; offense; crime.
	\end{compactitem}

% % ----------------------
% \subsection{Oxford English Dictionary} % *OED
% % ----------------------
	% \label{TRANSGRESSION-OED}

	% \begin{compactitem}
		% \item[]
	% \end{compactitem}

% ----------------------
\subsection{Scriptural Usage} % *SCRIPTURES
% ----------------------
	\label{TRANSGRESSION-SCRIPTURES}

	% % -----
	% \subsubsection{Old Testament} % *OT
	% % -----
		% \label{TRANSGRESSION-OT}
	
		% % --
		% \paragraph{KJV}
		% % --
			% \label{TRANSGRESSION-OT-KJV}
	
			% \begin{tabular}{ll}
				% Total occurrences in the OT & 
			% \end{tabular}
			
			% \begin{compactdesc}
				% \item[]
			% \end{compactdesc}
			
		% % --
		% \paragraph{Hebrew Bible}
		% % --
			% \label{TRANSGRESSION-HEBREW}
		
			% %
			% \subparagraph{\texorpdfstring{\he{sample word}}{}}
			% %
				% \label{PAGA} % whatever is in step bible without periods
			
				% \begin{compactdesc}
					% \item[HALOT , PDF ] \textbf{}
						% \begin{compactitem}
							% \item[1]
						% \end{compactitem}
					% \item[BDB , PDF ] \textbf{}
						% \begin{compactitem}
							% \item[1]
						% \end{compactitem}
					% \item[DCH PDF ] \textbf{}
						% \begin{compactitem}
							% \item[1]
						% \end{compactitem}
				% \end{compactdesc}
				
				% \begin{compactdesc}
					% \item[Some OT uses] \textbf{}
						% \begin{compactdesc}
							% \item[]
						% \end{compactdesc}
				% \end{compactdesc}
			
		% % --
		% \paragraph{Septuagint}
		% % --
			% \label{TRANSGRESSION-LXX}
		
			% %
			% \subparagraph{\texorpdfstring{\gr{sample word}}{}}
			% %
		
	% -----
	\subsubsection{New Testament} % *NT
	% -----
		\label{TRANSGRESSION-NT}
	
		% --
		\paragraph{KJV}
		% --
			\label{TRANSGRESSION-NT-KJV}
			
	
			\begin{tabular}{ll}
				Total occurrences in the NT & 8
			\end{tabular}
			
			\begin{compactdesc}
				\item[{\hyperref[ROMANS-4-15]{Romans 4:15}}] \phantom{.}
					\begin{compactitem}
						\item See 1 John 3:4 below.
					\end{compactitem} 
				\item[{\hyperref[1JOHN-3-4]{1 John 3:4}}] \phantom{.}
					\begin{compactitem}
						\item See the Greek here for sure.
						\item See also Romans 4:15 above.
					\end{compactitem}
			\end{compactdesc}
			
		% --
		\paragraph{Greek New Testament}
		% --
			\label{TRANSGRESSION-GREEK}
		
			%
			\subparagraph{\texorpdfstring{\gr{παράβασις}}{}}
			%
				\label{PARABASIS}
				
				See also \hyperref[HAMARTIA]{\grl{ἁμαρτία}}, \hyperref[ANOMIA]{\grl{ἀνομία}}, \hyperref[ADIKIA]{\grl{ἀδικία}}, 
			
				\begin{compactdesc}
					\item[BDAG 674a] \textbf{}
						\begin{compactitem}
							\item act of deviating from an established boundary or norm, \textit{overstepping, transgression} w. objective genitive
						\end{compactitem}
					\item[LSJ 1305a, PDF 1376] \textbf{}
						\begin{compactitem}
							\item[I.1] \textit{going aside, escape, deviation}
							\item[I.2] of the action of \textit{walking}
							\item[I.3] \textit{transition, passage}
							\item[II] \textit{overstepping, transgression, error, illusion}
							\item[III] \textit{parabasis}, a part of the Old Comedy, \textit{in which the Chorus came forward} and addressed the audience in the Poet's name
						\end{compactitem}
				\end{compactdesc}
				
				\begin{compactdesc}
					\item[Some NT uses] \textbf{}
						\begin{compactdesc}
							\item[Romans 2:23] \gr{ὃς ἐν νόμῳ καυχᾶσαι, διὰ τῆς παραβάσεως τοῦ νόμου τὸν θεὸν ἀτιμάζεις;}
							\item[Romans 5:14] \gr{ἀλλὰ ἐβασίλευσεν ὁ θάνατος ἀπὸ Ἀδὰμ μέχρι Μωϋσέως καὶ ἐπὶ τοὺς μὴ ἁμαρτήσαντας ἐπὶ τῷ ὁμοιώματι τῆς παραβάσεως Ἀδάμ, ὅς ἐστιν τύπος τοῦ μέλλοντος.}
						\end{compactdesc}
				\end{compactdesc}
		
	% -----
	\subsubsection{Book of Mormon} % *BOM
	% -----
		\label{TRANSGRESSION-BOM}
	
		\begin{tabular}{ll}
			Total occurrences in the BOM & 65
		\end{tabular}
		
		\begin{compactdesc}
			\item[{\hyperref[2NEPHI-9-6]{2 Nephi 9:6}}]
		\end{compactdesc}
		
	% -----
	\subsubsection{Doctrine and Covenants} % *DC
	% -----
		\label{TRANSGRESSION-DC}
	
		% \begin{tabular}{ll}
			% Total occurrences in the DC & 
		% \end{tabular}
		
		\begin{compactdesc}
			\item[{\hyperref[DC29]{DC 29}}] (In his original heading to this section in RB1, John Whitmer wrote, ``they seeing somewhat different upon the death of Adam (that is his transgression) therefor they made it a subject of Prayer \& enquired of the Lord.'')
		\end{compactdesc}
		
	% % -----
	% \subsubsection{Pearl of Great Price} % *PGP
	% % -----
		% \label{TRANSGRESSION-PGP}
	
		% \begin{tabular}{ll}
			% Total occurrences in the PGP & 
		% \end{tabular}
		
		% \begin{compactdesc}
			% \item[]
		% \end{compactdesc}
		
% ----------------------
\subsection{Key Phrases} % *KEYPHRASES *PHRASES
% ----------------------
	\label{TRANSGRESSION-PHRASES}

	\begin{compactdesc}
		\item[transgression of Adam] Job 31:33; Romans 5:14; 2 Nephi 2:22; Mosiah 3:11; Alma 22:12; DC 29:40; Moses 5:10, 6:53; Articles of Faith 2
			\begin{compactdesc}
				\item[Bible anchor verses] Romans 5:14
					\begin{compactdesc}
						\item[Romans 5:14] \gr{τῆς παραβάσεως Ἀδάμ}
					\end{compactdesc}
				\item[Commentary] \phantom{.}
					\begin{compactitem}
						\item The BDAG definition appropriate for Romans 5:14 appears to be an ``act of deviating from an established boundary or norm.'' The fact that it is referred to in Romans 5:14 as a \gr{παράβασις} and not a \gr{ἁμαρτία} in connection with Adam is probably very important; those are not the same thing.
						\item Note John Whitmer's original heading to \hyperref[DC29]{DC 29}, mentioned above, where he talks about ``the death of Adam (that is his transgression).''
					\end{compactitem}
			\end{compactdesc}
	\end{compactdesc}
	
% ----------------------
\subsection{Bible Dictionaries} % *DICTIONARIES
% ----------------------
	\label{TRANSGRESSION-BD}
	
		% -----
		\subsubsection{Theological Dictionary of the New Testament Abridged}
		% -----
			\label{TRANSGRESSION-BD-TDNTA}
		
			% --
			\paragraph{\texorpdfstring{\gr{παράβασις}}{}}
			% --
			
				This word means ``striding to and fro,'' ``stepping over,'' ``transgressing,'' ``violating'' (cf. in the LXX Ps. 101:3). In the NT it denotes sin in relation to the law. In Rom. 2:23 the Jew dishonors God by transgressing the law, and \textbf{in 4:15 the law brings wrath because there is transgression only where there is law. Between Adam and Moses there is sin but no \gr{παράβασις} because the law is not yet given.} In Gal. 3:19 the law is given to show that evil deeds are transgressions of God's will. In 1 Tim. 2:14 Eve is deceived first and becomes a transgressor by violating God's prohibition. In Heb. 2:2 every transgression of the law carries a due penalty, and in 9:15 Christ's death serves to remit the transgressions committed under the old covenant, i.e., the acts of disobedience against God's law that have brought Israel into guilt against God. 

% % ----------------------
% \subsection{Restoration Usage} % *RESTORATION
% % ----------------------
	% \label{TRANSGRESSION-RESTORATION}

	% % -----
	% \subsubsection{Joseph Smith} % *JS
	% % -----
		% \label{TRANSGRESSION-JS}
	
		% \begin{compactdesc}
			% \item[DATE/SOURCE]
		% \end{compactdesc}
	
	% % -----
	% \subsubsection{Brigham Young} % *BY
	% % -----
		% \label{TRANSGRESSION-BY}
	
		% \begin{compactdesc}
			% \item[DATE/SOURCE]
		% \end{compactdesc}
	
% ----------------------
\subsection{Commentary and Studies} % *COMMENTARY *STUDIES
% ----------------------
	\label{TRANSGRESSION-COMMENTARY}

	% -----
	\subsubsection{Key Points}
	% -----
		\label{TRANSGRESSION-KEYS}
	
		\begin{compactitem}
			\item There is transgression only where there is law (see \hyperref[TRANSGRESSION-BD-TDNTA]{here}).
			\item Sin and transgression come apart:
				\begin{compactitem}
					\item You can have sin without transgression, in cases where there is no law.
					\item You can have sin and transgression, in cases where there is a law and someone has violated it and sinned.
					\item You can have transgression without sin, which may be the case for Adam, and is the case for us, in times when we violate the law but don't actually sin.
				\end{compactitem}
			\item There is a connection or conceptual equivalence relation between transgression and death---in some places, these concepts may be substitutable for each other (see the early preface to DC 29 \hyperref[EMS-1830-JS-SEP-1830-A]{here}).
		\end{compactitem}
		
	% -----
	\subsubsection{The difference between sin and transgression}
	% -----
		
		See \hyperref[SIN-sin-transgression]{these notes} at the entry for \hyperref[SIN]{Sin}.